% ============================================================
%  Глава 6 — Что такое регулятор: от P к ЛКР
% ============================================================
\chapter{Что такое регулятор}
\label{ch:controllers}

\section{Вопрос}

У нас есть модель $\vect{x}_{k+1} = \mat{A}_d\vect{x}_k + \mat{B}_d\vect{u}_k$.
Робот «понимает», как состояние реагирует на команду. \textbf{Какую
команду $\vect{u}_k$ ему подавать на каждом тике}, чтобы достичь цели?

Эта задача --- задача синтеза \emph{регулятора}.

\begin{ideabox}
Регулятор --- это \emph{функция}, которая по текущему состоянию (и,
возможно, желаемому состоянию) выдаёт управление:
\[
  \vect{u}_k = \pi(\vect{x}_k, \vect{x}^\text{ref}_k).
\]
Слово «регулятор» --- из ТАУ. Часто говорят также «контроллер»,
«controller», «policy». Это одно и то же.
\end{ideabox}

Идея этой главы --- \emph{интуитивно} объяснить, чем отличаются простой
П-регулятор, PID и ЛКР. К концу главы будет понятно, \emph{зачем} мы вообще
идём в матрицы $\mat{Q}, \mat{R}$ и почему нельзя обойтись пропорциональным
коэффициентом.

\section{П-регулятор (proportional)}

Самый простой регулятор: чем дальше от цели --- тем сильнее тяни.

\begin{equation}
  u_k = -k_p \cdot (x_k - x^\text{ref}_k).
  \label{eq:P_controller}
\end{equation}

Здесь $k_p > 0$ --- \emph{коэффициент усиления} (gain). Чем он больше,
тем «жёстче» регулятор. Скаляр; в многомерном случае --- вектор/матрица.

\begin{examplebox}[title={Скаляр, удержание скорости}]
Хотим, чтобы робот ехал ровно $v_\text{target}$. Текущая скорость $v$.
Тогда $u = -k_p (v - v_\text{target})$. Если $v < v_\text{target}$ ---
команда положительная, разгон. Если $v > v_\text{target}$ --- отрицательная,
торможение.
\end{examplebox}

\textbf{Проблемы П-регулятора:}

\begin{itemize}
  \item \textbf{Установившаяся ошибка.} Если есть сопротивление (трение,
        уклон), при достижении $v = v_\text{target}$ команда обнуляется,
        но трение остаётся --- робот замедлится. На равновесии останется
        ошибка $\ne 0$.
  \item \textbf{Колебания.} Большой $k_p$ $\to$ перерегулирование $\to$
        качается вокруг цели.
  \item \textbf{Не учитывает динамику.} П-регулятор «слепой»: ему всё равно,
        что моторы инерционны. Если $\tau_v$ велико --- регулятор
        будет «опаздывать» на каждом тике.
\end{itemize}

\section{PID}

Классическая попытка починить проблемы П: добавить \emph{интеграл} ошибки
(чтобы убрать установившуюся) и \emph{производную} (чтобы предугадать).

\begin{equation}
  u_k = -k_p \cdot e_k - k_i \cdot \sum_{j=0}^{k} e_j \Delta t - k_d \cdot (e_k - e_{k-1})/\Delta t.
  \label{eq:PID_controller}
\end{equation}

PID --- старший товарищ. Работает на всём: от лифта до автопилота. Но
у него тоже минусы:

\begin{itemize}
  \item Подбор $k_p, k_i, k_d$ --- искусство (Зиглер--Никольс, ручное
        крутение). На многомерных задачах перестаёт быть удобным.
  \item Не оптимален по никакому критерию: «работает» $\ne$ «работает
        лучше всего».
  \item \textbf{Не знает модель} объекта. Если бы знал $\mat{A}_d, \mat{B}_d$
        --- мог бы предсказывать, а не реагировать. PID не предсказывает.
\end{itemize}

\section{Линейно-квадратичный регулятор (ЛКР)}

ЛКР (linear-quadratic regulator, LQR) --- иначе подходит к задаче.

\begin{ideabox}
\textbf{Идея ЛКР:} вместо ручного подбора коэффициентов
\emph{сформулировать} цель как минимизацию специальной функции
\textbf{критерия качества}, и пусть математика сама найдёт лучший
регулятор для этой цели.
\end{ideabox}

Критерий --- сумма «штрафов» по всему будущему:

\begin{formulabox}[title={ЛКР-критерий качества}]
\begin{equation}
  J = \sum_{k=0}^{\infty} \Big(\vect{x}_k\T \mat{Q}\,\vect{x}_k +
                                  \vect{u}_k\T \mat{R}\,\vect{u}_k\Big).
  \label{eq:lqr_cost}
\end{equation}
\end{formulabox}

Здесь:

\begin{itemize}
  \item $\mat{Q} \in \Real^{n\times n}$ --- штраф за \textbf{отклонения
        состояния} от нуля (предполагаем, что цель --- $\vect{x} = 0$;
        если цель ненулевая, перейдём в координаты ошибки).
  \item $\mat{R} \in \Real^{r\times r}$ --- штраф за \textbf{величину
        управления}. Большой $\mat{R}$ --- регулятор будет «экономить»
        моторы, малый --- «топить педаль в пол».
  \item Обе матрицы --- симметричные, положительно (полу-)определённые.
        На практике диагональные:
        $\mat{Q} = \diag(q_1, \ldots, q_n)$, $\mat{R} = \diag(r_1, \ldots, r_r)$.
\end{itemize}

\subsection{Что значит «положительно определённая»}

\begin{notebox}
Матрица $\mat{Q}$ \textbf{положительно (полу-)определённая}, если для
любого ненулевого $\vect{x}$ выполняется $\vect{x}\T \mat{Q}\,\vect{x} > 0$
(или $\ge 0$ --- полу-).

Для диагональной $\mat{Q} = \diag(q_1, \ldots, q_n)$ это сводится к
$q_i \ge 0$ (или строгая положительность для всех $q_i$, тогда уже
positive definite).

Зачем требование: штраф $\vect{x}\T \mat{Q}\,\vect{x}$ должен быть
\emph{неотрицательным}, иначе оптимизатор «разгонит» состояние до
$-\infty$ и получит бесконечную выгоду. Аналогично $\mat{R}$ должна
быть строго положительно определённая --- иначе можно бесконечно
увеличивать $\vect{u}$ без штрафа.
\end{notebox}

\subsection{Зачем квадратичный штраф?}

Почему именно $\vect{x}\T\mat{Q}\vect{x}$, а не, скажем, $|\vect{x}|$?

\begin{itemize}
  \item \textbf{Дифференцируем} $\to$ можно решить аналитически.
  \item Маленькие ошибки штрафуются \emph{ещё меньше} (квадрат малого числа
        --- ещё меньше), большие --- \emph{гораздо больше} (квадрат
        $5 = 25$, квадрат $10 = 100$). Регулятор сильно реагирует на
        большие отклонения, мягко --- на мелкие.
  \item Удобная аналогия с энергией / расстоянием. Норма $\|\vect{x}\|^2
        = \vect{x}\T\vect{x}$, а $\vect{x}\T\mat{Q}\vect{x}$ --- «взвешенная»
        норма.
\end{itemize}

\subsection{Результат ЛКР}

\textbf{Ключевой факт.} Если объект линейный, дискретный, управляемый,
и критерий \eqref{eq:lqr_cost} с положительно определёнными $\mat{Q}, \mat{R}$,
то оптимальный регулятор --- \textbf{линейная обратная связь}:

\begin{formulabox}[title={ЛКР-закон управления}]
\begin{equation}
  \vect{u}_k = -\mat{K}\,\vect{x}_k,
  \label{eq:lqr_law}
\end{equation}
\end{formulabox}

где $\mat{K}$ --- матрица $r \times n$, не зависящая от $k$ (постоянная).
Она вычисляется \emph{один раз}, оффлайн, через \emph{уравнение Риккати}.
\emph{Как именно} --- следующая глава.

\section{Интуитивная карта}

\begin{center}
\begin{tabular}{lp{8cm}}
\toprule
\textbf{Регулятор} & \textbf{Идея}\\
\midrule
П    & «Тяни сильнее, если дальше».\\
PID  & П + интеграл + производная (борьба с трением и инерцией).\\
ЛКР  & Минимизируй квадратичный критерий $J = \sum (\vect{x}\T \mat{Q}\vect{x}
       + \vect{u}\T \mat{R}\vect{u})$. Решение --- $\vect{u} = -\mat{K}\vect{x}$.
       \textbf{Оптимален по выбранному критерию.}\\
MPC  & Каждый тик решай ЛКР-подобную задачу \emph{на короткий горизонт
       вперёд}, с \emph{учётом ограничений} $\vect{u}_\text{min} \le \vect{u} \le \vect{u}_\text{max}$.\\
\bottomrule
\end{tabular}
\end{center}

\begin{ideabox}
\textbf{Логическая цепочка:}
\begin{itemize}
  \item П --- слепая реакция.
  \item PID --- реакция с памятью и предсказанием в эвристике.
  \item ЛКР --- \emph{оптимальная} реакция при заданном критерии и
        \emph{знании модели}. Не учитывает ограничения.
  \item MPC --- ЛКР с учётом ограничений на $\vect{u}$ и предсказанием
        на конечный горизонт $N$.
\end{itemize}
Каждый следующий --- более мощный и при этом более вычислительно дорогой.
В проекте используются ЛКР и MPC одновременно (MPC при \texttt{solver=clip}
сводится к ЛКР с клиппингом --- эффективный компромисс).
\end{ideabox}

\section{Связь с проектом}

\begin{codebox}[title={В проекте есть оба:}]
\begin{itemize}
  \item \filepath{pi\_nodes/control/lqr\_controller.py} --- чистый ЛКР
        (используется как baseline и для сравнения).
  \item \filepath{pi\_nodes/control/mpc\_controller.py} --- MPC. В режиме
        \texttt{solver = "clip"} он \emph{буквально} вычисляет ЛКР-усиление
        $\mat{K}_\text{first}$ (первая строка лифтированной QP-задачи),
        применяет его и клиппит в ограничения.
\end{itemize}
\end{codebox}

В следующих двух главах разберём ЛКР, потом MPC.
